% A worked example of the cktimg LaTeX package.
%
% Compile from the repository root, so the relative netlist paths resolve:
%
%   zig build -Dlatex_renderer=true              # builds zig-out/bin/cktimg-tex
%   TEXINPUTS=latex: pdflatex -shell-escape examples/latex/figure.tex
%
% Without -shell-escape, generate the figures first and compile normally:
%
%   ./zig-out/bin/cktimg-tex tests/fixtures/current_mirror.spice current_mirror.cktimg.tex
%   ./zig-out/bin/cktimg-tex tests/fixtures/two_stage_miller.spice two_stage_miller.cktimg.tex
%   TEXINPUTS=latex: pdflatex examples/latex/figure.tex
%
% The package figures out which of those two you did; see latex/cktimg.sty.

\documentclass[11pt]{article}
\usepackage[margin=1in]{geometry}

% [auto] is the default: shell escape when available, a pre-generated .tex otherwise.
% Use [pregenerated] on Overleaf, or [shell] to make a missing -shell-escape a hard
% error instead of a confusing "file not found".
\usepackage{cktimg}

% cktimg-tex is not on PATH when working inside the repository, so point at the build
% output. Delete this line once the binary is installed system-wide.
\renewcommand{\cktimgbinary}{./zig-out/bin/cktimg-tex}

\begin{document}

\section*{Schematics from netlists}

Each figure below is drawn straight from the SPICE deck named in the
\verb|\cktimg| call --- there is no intermediate drawing step and no coordinates in
this document. Placement, orientation and wire routing are decided by cktImg; the
output is ordinary TikZ, so it scales, prints and embeds like anything else in the
paper.

\begin{figure}[h]
  \centering
  \cktimg{tests/fixtures/current_mirror.spice}
  \caption{A current mirror. The reference leg sets the current the output leg copies.}
\end{figure}

\begin{figure}[h]
  \centering
  % The optional argument is handed to cktimg-tex verbatim, so any flag the program
  % accepts works here. A lint.zon can widen the grid, recolour the wires, or spend
  % more search effort on the column order:
  %
  %   \cktimg[--config lint.zon]{tests/fixtures/two_stage_miller.spice}
  \cktimg{tests/fixtures/two_stage_miller.spice}
  \caption{A two-stage Miller-compensated amplifier.}
\end{figure}

\end{document}
